\section{Introduction}

The rapid integration of Large Language Models (LLMs) into the software engineering lifecycle has led to the proliferation of autonomous coding agents, such as Claude Code and OpenCode \cite{mcp2025}. While these agents demonstrate high proficiency in isolated, single-session tasks, their practical utility diminishes when operating in complex, multi-day, or multi-developer environments. The primary failure mode in this paradigm is context amnesia: agents function as ephemeral silos, losing vital architectural rationale, constraints, and state transitions across session boundaries.

Current approaches attempt to resolve this discontinuity by treating project state as a retrieval problem, utilizing vector databases and semantic search (e.g., \cite{honcho2026, mem02026}) to inject historical transcripts into an agent's context window. However, this ``Project Memory'' approach suffers from a fundamental flaw: it conflates unverified agent proposals with authoritative project decisions. Over time, memory stores accumulate noise, obsolete hypotheses, and discarded code paths, leading to cascading hallucinations and compounding rework. 

We hypothesize that resolving multi-agent discontinuity requires a shift from passive \textit{Memory} to formal \textit{Project Truth}. To preserve accountability and continuity across agentic handoffs \cite{a2a2025}, systems must reconcile disparate claims against deterministic evidence, such as version control commits and test suite outcomes, establishing a governed lifecycle of commitments \cite{singh1999commit}.

In this paper, we present the \textit{Bridge} framework, an authoritative project truth and reconciliation substrate designed to facilitate interoperable software engineering agents. Through a sequence of controlled experiments (\pendingdata{EXP-001 and EXP-002 results}), we evaluate the impact of structured work-state transfer and persistent project intelligence on agent performance.

Our core contributions are:
\begin{itemize}
    \item A formalization of the agent interoperability boundary, transitioning from context transfer to governed project truth.
    \item An empirical evaluation of structured work-state handoffs across heterogeneous agents (\pendingdata{reference to EXP-001 empirical claims}).
    \item A conceptual framework for state reconciliation and commitment control in multi-agent software engineering.
\end{itemize}
